\SourcePage{653}
\SGAsection{4}{关于 $\operatorname{Pic}^\tau_{X/S}$ 的有限性定理}
\textbf{4.1。}设 $S$ 为\Term{scheme}{概形}，$G$ 为 $S$ 上的\Term{abelian-group}{交换群}\Term{functor}{函子}。以 $G^0$（分别为 $G^\tau$）表示如下群子\Term{functor}{函子}：其 $T$ 值点为那些 $f:T\to G$，使得对每个\Term{algebraically-closed-field}{代数闭域} $K$ 及 $T$ 的每个 $K$ 值点 $t$，存在\Term{connected}{连通}的 $K$-\Term{scheme}{概形} $Z$、$Z$ 的两个 $K$ 值点 $z,z_0$ 和\Term{morphism}{态射} $g:Z\to G$，满足 $g(z)=f(t)$、$g(z_0)=0\in G(K)$（分别为：存在非零整数 $n$，使 $n f(t)\in G^0(K)$）。

显然，$G^0$（分别为 $G^\tau$）的构造与任意\Term{base-change}{基变换}交换；群\Term{functor}{函子}的\Term{morphism}{态射} $G\to G'$ 诱导\Term{morphism}{态射} $G^0\to G'^0$（分别为 $G^\tau\to G'^\tau$）。若 $S$ 为域，且 $G$ 由\Term{locally-finite-type}{局部有限型}群\Term{scheme}{概形}表示，则 $G^0$ 由\Term{identity-element}{单位元}的\Term{connected-component}{连通分支}表示；它是\Term{finite-type}{有限型}、\Term{geometrically-irreducible}{几何不可约}的开（也是闭）子\Term{scheme}{概形}（SGA 3，第 VI$_A$ 讲 2.4）。

\begin{SGAresult}{\Term{lemma}{引理}}{4.2}
设 $\Phi:G\to G'$ 是\Term{scheme}{概形} $S$ 上两个\Term{abelian-group}{交换群}\Term{functor}{函子}之间的\Term{morphism}{态射}。假设对每个\Term{algebraically-closed-field}{代数闭域} $K$ 及每个\Term{morphism}{态射} $\operatorname{Spec}(K)\to S$，$G_K$（分别为 $G'_K$）均可由\Term{locally-finite-type}{局部有限型}群\Term{scheme}{概形}表示，且 $\Phi_K$ 可由\Term{finite-type}{有限型}（即\Term{quasi-compact}{拟紧}）\Term{morphism}{态射}表示。则
\[\Phi^{-1}(G'^\tau)=G^\tau.\]
\end{SGAresult}
\SourcePage{654}
\begin{proof}
我们已指出 $\Phi(G^\tau)\subset G'^\tau$。反向包含关系由下述事实得到：$g\in G(K)$ 属于 $G^\tau(K)$，当且仅当序列 $ng$（$n\in\mathbf N$）\Term{quasi-compact}{拟紧}；由于 $\Phi$ \Term{quasi-compact}{拟紧}，只须序列 $\Phi(ng)=n\Phi(g)$ 具有同一性质即可。
\end{proof}
\begin{SGAresult}{\Term{lemma}{引理}}{4.3}
设 $X,Y$ 均\Term{proper}{固有}于\Term{scheme}{概形} $S$，$f:X\to Y$ 为\Term{surjective}{满射}\Term{morphism}{态射}（分别为\Term{ample}{丰沛}\Term{sheaf}{层} $\mathscr O_Y(1)$ 的一个\Term{section}{截面} $\varphi$ 的零点\Term{scheme}{概形}的包含\Term{morphism}{态射}）。则
\[(f^*)^{-1}\operatorname{Pic}^\tau_{X/S}=\operatorname{Pic}^\tau_{Y/S}.\]
\end{SGAresult}
\begin{proof}
这是 4.2、第 XII 讲 1.2、3.1 及 3.5（分别为 3.8）的\Term{form}{形式}推论。
\end{proof}
\textbf{4.4。}设 $k$ 为\Term{algebraically-closed-field}{代数闭域}，$X$ 为\Term{proper}{固有} $k$-\Term{scheme}{概形}，$L$ 为 $X$ 上的\Term{invertible-sheaf}{可逆层}。称 $L$ \Term{algebraic-equivalence}{代数等价}于 $0$（分别为 $\tau$-\Term{tau-equivalence}{等价}于 $0$），若存在\Term{connected}{连通}\Term{finite-type}{有限型} $k$-\Term{scheme}{概形} $T$、$X_T=X\times_kT$ 上的\Term{invertible-sheaf}{可逆层} $M$，以及 $T$ 的两个 $k$ 值点 $t,t_0$，使 $L\simeq M_t$、$\mathscr O_X\simeq M_{t_0}$（分别为：存在非零整数 $n$，使 $L^{\otimes n}$ \Term{algebraic-equivalence}{代数等价}于 $0$）。容易看出，这也\Term{equivalence}{等价}于：$L$ 对应的 $\operatorname{Pic}_{X/k}$ 的 $k$ 值点 $\lambda$ 属于 $\operatorname{Pic}^0_{X/k}(k)$（分别为 $\operatorname{Pic}^\tau_{X/k}(k)$）。

还要回顾，称 $L$ \Term{numerical-equivalence}{数值等价}于 $0$，若对 $X$ 的每条\Term{integral}{整}闭\Term{curve}{曲线} $Y$，\Term{intersection-number}{交数} $\langle c_1(L).Y\rangle$ 均为 $0$（见 2.1 或第 X 讲 4.5）。
\SourcePage{655}
\begin{SGAresult}{命题}{4.5（\Term{functorial}{函子性}）}
设 $f:X'\to X$ 为两个\Term{proper}{固有} $k$-\Term{scheme}{概形}之间的\Term{morphism}{态射}，$L$ 为 $X$ 上的\Term{invertible-sheaf}{可逆模}。则：
\begin{enumerate}
\item[(i)] 若 $L$ $\tau$-\Term{tau-equivalence}{等价}于 $0$（分别为\Term{numerical-equivalence}{数值等价}于 $0$），则 $f^*L$ 亦然。
\item[(ii)] 反之，若 $f$ \Term{surjective}{满射}，且 $f^*L$ $\tau$-\Term{tau-equivalence}{等价}于 $0$（分别为\Term{numerical-equivalence}{数值等价}于 $0$），则 $L$ 亦然。
\end{enumerate}
\end{SGAresult}
\begin{proof}
对于 $\tau$-\Term{tau-equivalence}{等价}，断言 (i)（分别为 (ii)）立即由 4.1（分别为 4.3）得到。另一方面，设 $Y'$ 为 $X'$ 的一条\Term{integral}{整}闭\Term{curve}{曲线}，$Y=f(Y')$ 为其像，赋予诱导的\Term{reduced}{约化}结构。若 $\dim Y=1$，则有\Term{projection-formula}{投影公式}（第 X 讲 4.4.4）
\[\langle c_1(f^*L).Y'\rangle=[k(Y'):k(Y)]\langle c_1(L).Y\rangle.\tag{4.5.1}\]
这里 $k(Y')$（分别为 $k(Y)$）是 $Y'$（分别为 $Y$）的\Term{function-field}{函数域}。若 $\dim Y=0$，则 $(f^*L)_{Y'}\simeq\mathscr O_{Y'}$。因此，若 $L$ \Term{numerical-equivalence}{数值等价}于 $0$，便有 $\langle c_1(f^*L).Y'\rangle=0$，这就证明了 (i)。

反之，若 $f$ \Term{surjective}{满射}，则对 $X$ 的每条\Term{integral}{整}闭\Term{curve}{曲线} $Y$，都可找到 $X'$ 的一条\Term{integral}{整}闭\Term{curve}{曲线} $Y'$ 使 $Y=f(Y')$：取 $Y$ 的\Term{generic-point}{泛点} $y$、\Term{fiber}{纤维} $f^{-1}(y)$ 中在 $k(y)$ 上\Term{algebra}{代数}的点 $y'$，并令 $Y'$ 为 $y'$ 的\Term{closure}{闭包}。于是 4.5.1 表明 $\langle c_1(f^*L).Y'\rangle=0$ 蕴含 $\langle c_1(L).Y\rangle=0$，这就证明了 (ii)。
\end{proof}
\begin{SGAresult}{\Term{theorem}{定理}}{4.6（$\operatorname{Pic}^\tau_{X/k}$ 的数值刻画）}
设 $k$ 为\Term{algebraically-closed-field}{代数闭域}，$X$ 为\Term{proper}{固有} $k$-\Term{scheme}{概形}，$L$ 为 $X$ 上的\Term{invertible-sheaf}{可逆层}。则下列条件\Term{equivalence}{等价}：
\begin{enumerate}
\item[(i)] $L$ $\tau$-\Term{tau-equivalence}{等价}于 $0$。
\item[(ii)] 对 $X$ 上的每个\Term{coherent}{凝聚}\Term{sheaf}{层} $F$，有
\SourcePage{656}
\[\chi(F\otimes L)=\chi(F).\]
\item[(iii)] $L$ \Term{numerical-equivalence}{数值等价}于 $0$。
\end{enumerate}
若 $X/k$ \Term{projective}{射影}，并配有\Term{ample}{丰沛}\Term{invertible-sheaf}{可逆模} $H=\mathscr O_X(1)$，则以上条件还\Term{equivalence}{等价}于下列条件：
\begin{enumerate}
\item[(iv)] 对每个整数 $m$（正整数或负整数），$L^{\otimes m}(1)$ 均\Term{ample}{丰沛}。
\item[(v)] （若 $X$ \Term{irreducible}{不可约}。）对每对整数 $m,n$，有
\[\chi(L^{\otimes m}(n))=\chi(\mathscr O_X(n)).\]
\item[(vi)] （若 $X$ \Term{irreducible}{不可约}且\Term{dimension}{维数}为 $r\geq2$。）有
\[\langle c_1(L).c_1(H)^{r-1}\rangle=0,\qquad
\langle c_1(L)^2.c_1(H)^{r-2}\rangle=0.\]
\end{enumerate}
\end{SGAresult}
\begin{proof}
先证 (i) 蕴含 (ii)。应用\Term{resolution}{分解}\Term{lemma}{引理}（EGA III 3.1）及 4.5 (i)，可归结到 $X$ \Term{integral}{整}且 $F=\mathscr O_X$ 的情形。此时须证明\Term{polynomial}{多项式} $\chi(L^{\otimes m})$ 为常数。为此可将 $L$ 换为它的某个倍数，因而可设 $L$ \Term{algebraic-equivalence}{代数等价}于 $0$。此时，若 $L$ 经由 $X\times_kT$ 上的\Term{invertible-sheaf}{可逆层} $M$ 所定义的\Term{algebra}{代数}族变形为 $0$，其中 $T$ \Term{connected}{连通}且\Term{finite-type}{有限型}，则\Term{polynomial}{多项式} $\chi(M_t^{\otimes m})$ 与 $t\in T$ 无关（EGA III 7.9.5）；所求蕴含关系由此得到。

设 $Y$ 为 $X$ 的一条\Term{integral}{整}闭\Term{curve}{曲线}。在 (ii) 中取 $F=\mathscr O_Y$，得到 $\chi(\mathscr O_Y)=\chi(L_Y)$。于是公式（“Riemann \Term{theorem}{定理}”）
\[\chi(L_Y)=\langle c_1(L).Y\rangle+\chi(\mathscr O_Y)\]
给出 $\langle c_1(L).Y\rangle=0$，故 (ii) 蕴含 (iii)。

取 \Term{chow-presentation}{Chow 表示} $f:X'\to X$，即 $X'$ \Term{projective}{射影}且 $f$ \Term{surjective}{满射}。若已知对 $X'$ 有 (iii) 蕴含 (i)，则由 4.5 对 $X$ 也有该蕴含关系。
\SourcePage{657}
因此以下设 $X$ \Term{projective}{射影}，并配有\Term{ample}{丰沛}\Term{sheaf}{层} $H$。

设 $X_i$ 为 $X$ 的各\Term{irreducible-component}{不可约分支}，赋予诱导的\Term{reduced}{约化}结构。对\Term{finite}{有限}\Term{surjective}{满射}\Term{morphism}{态射} $\coprod X_i\to X$ 应用 4.5 及 EGA III 2.6.2，便将 (iii) 蕴含 (i) 和 (iv) 的证明归结到 $X$ \Term{integral}{整}的情形。

若 $X$ \Term{irreducible}{不可约}，则显然 (ii) 蕴含 (v)，(iii) 蕴含 (vi)；结合 2.1，(v) 蕴含 (vi)，且 $X$ 的 (vi) \Term{equivalence}{等价}于 $X_{\mathrm{red}}$ 的 (vi)（第 X 讲 4.4.1）。

因而只须在 $X$ \Term{integral}{整}的假设下，证明 (vi) 蕴含 (i) 和 (iv)，以及 (iv) 蕴含 (i)。此时，(vi)（分别为 (iv)）说明各 $L^{\otimes m}$（$m\in\mathbf N$）的类所成的族 $\mathscr L$ 满足 3.13 (iii)（分别为 (iv)），故为\Term{bounded-family}{有界族}。因此存在 $X_T$ 上的\Term{invertible-sheaf}{可逆层} $\mathscr L$，其中 $T$ 在 $k$ 上\Term{finite-type}{有限型}，使之界定该族。取两个不同的整数 $p,q$，使 $L^{\otimes p}$ 和 $L^{\otimes q}$ 对应于 $T$ 的同一\Term{connected-component}{连通分支}；则 $L^{\otimes(p-q)}$ \Term{algebraic-equivalence}{代数等价}于 $0$。故 (vi)（分别为 (iv)）蕴含 (i)。

最后，同样由 (vi) 可得存在整数 $N$，使对每个整数 $m$，$L^{\otimes m}\otimes H^{\otimes N}$ 都\Term{ample}{丰沛}。于是对每个 $m$，$(L^{\otimes m}\otimes H)^{\otimes N}$ 都\Term{ample}{丰沛}，故 (vi) 蕴含 (iv)。证明完毕。
\end{proof}
\begin{SGAresult}{\Term{theorem}{定理}}{4.7}
设 $S$ 为\Term{noetherian}{诺特}\Term{scheme}{概形}，$X$ 为\Term{proper}{固有} $S$-\Term{scheme}{概形}。则：
\begin{enumerate}
\item[(i)] 包含\Term{morphism}{态射} $\operatorname{Pic}^\tau_{X/S}\to\operatorname{Pic}_{X/S}$ 可由\Term{open-immersion}{开浸入}表示。
\item[(ii)] 若 $X$ 在 $S$ 上\Term{projective}{射影}且\Term{integral}{整}，则 $\operatorname{Pic}^\tau_{X/S}\to\operatorname{Pic}_{X/S}$
\SourcePage{658}
也可由\Term{closed-immersion}{闭浸入}表示。
\item[(iii)] $\operatorname{Pic}^\tau_{X/S}$ 在 $S$ 上\Term{finite-type}{有限型}，即 $X/S$ 各\Term{fiber}{纤维}上 $\tau$-\Term{tau-equivalence}{等价}于 $0$ 的\Term{invertible-sheaf}{可逆层}的类所成的族为\Term{bounded-family}{有界族}。
\end{enumerate}
\end{SGAresult}
\begin{proof}
在 (i)（分别为 (ii)）中，须对任意 $S$-\Term{scheme}{概形} $T$ 的\Term{morphism}{态射} $T\to\operatorname{Pic}_{X/S}$，证明 $\operatorname{Pic}^\tau_{X/S}$ 的\Term{inverse-image}{逆像}可由 $T$ 的开（分别为\Term{open-and-closed}{开闭}）子\Term{scheme}{概形}表示。只须取位于 $S$ 的某个\Term{affine}{仿射}\Term{open-subset}{开集} $\operatorname{Spec}(A)$ 之上的 $T=\operatorname{Spec}(B)$。此时 $T$ 是\Term{affine}{仿射} $S$-\Term{scheme}{概形} $T_\alpha=\operatorname{Spec}(B_\alpha)$ 所成\Term{filtered}{滤过}\Term{inverse-system}{投射系}的\Term{limit}{极限}，其中 $B_\alpha$ 遍历 $B$ 的\Term{finite-type}{有限型} $A$ \Term{subalgebra}{子代数}。由 3.1，所讨论的\Term{morphism}{态射} 来自某个 $T_\alpha\to\operatorname{Pic}_{X/S}$。故可设 $T$ 在 $S$ 上\Term{finite-type}{有限型}。

\Term{morphism}{态射} $T\to\operatorname{Pic}_{X/S}$ 对应于一个 f.p.p.f. \Term{morphism}{态射} $R\to T$ 和 $X_R$ 上的\Term{invertible-sheaf}{可逆层} $L$。由 f.p.p.f. \Term{descent}{下降}，可将 $T$ 换为 $R$。作\Term{base-change}{基变换} $R\to S$，并注意 $\operatorname{Pic}^\tau_{X/S}$ 的构造与任意\Term{base-change}{基变换}相容（4.1），则 (i)、(ii) 归结为研究由 $X$ 上的\Term{invertible-sheaf}{可逆层} $L$ 定义的\Term{section}{截面} $S\to\operatorname{Pic}_{X/S}$。

须证明由使 $L_s$ $\tau$-\Term{tau-equivalence}{等价}于 $0$ 的 $s\in S$ 所成的\Term{set}{集合} $S_0$ 是开的，并在 (ii) 的条件下也是闭的。由于 $S$ \Term{noetherian}{诺特}，这相当于对 $s\in S_0$ 证明以下三点：
\begin{enumerate}
\item[(a)] $s$ 的每个\Term{generalization}{泛化}仍在 $S_0$ 中；
\item[(b)] $s$ 附近的每个\Term{specialization}{特化}仍在 $S_0$ 中；
\item[(c)] 在 (ii) 的条件下，$s$ 的每个\Term{specialization}{特化}均在 $S_0$ 中。
\end{enumerate}
\SourcePage{659}
取 \Term{chow-presentation}{Chow 表示} $f:X'\to X$，其中 $X'$ 在 $S$ 上\Term{projective}{射影}，$f$ \Term{surjective}{满射}。若对 $X'$ 已知 (a)、(b)、(c) 及 (iii)，则由 4.3 和 3.5 对 $X$ 也成立。因此可设 $X$ 在 $S$ 上\Term{projective}{射影}。

取 $X/S$ 的\Term{ample}{丰沛}\Term{sheaf}{层} $H$。若对某个 $s\in S$，$(L^{\otimes m}\otimes H)_s$ \Term{ample}{丰沛}，则对 $s$ 的每个\Term{generalization}{泛化} $t$，$(L^{\otimes m}\otimes H)_t$ 都\Term{ample}{丰沛}（EGA III 4.7.1）。由 4.6 (i) 与 (iv) 的\Term{equivalence}{等价}性，便得 (a)。

为证明 (b)、(c) 及 4.7 (iii)，可设 $S$ \Term{integral}{整}，且在 (b)、(c) 中 $s$ 为\Term{generic-point}{泛点}。（对 (b)、(c) 这是显然的；对 (iii)，则由 3.4 (i) 应用于\Term{morphism}{态射} $\coprod S_i\to S$ 得到，其中 $S_i$ 为 $S$ 的各\Term{irreducible-component}{不可约分支}，赋予诱导的\Term{reduced}{约化}结构。）

存在\Term{integral}{整}\Term{scheme}{概形} $S'$ 及 f.p.p.f. \Term{morphism}{态射} $S'\to S$，使 $X_{S'}$ 的每个\Term{irreducible-component}{不可约分支} $X'_i$ 赋予诱导的\Term{reduced}{约化}结构后，在 $S'$ 上\Term{integral}{整}（第 XII 讲 1.3）。为证明 (b)、(iii)，只须对 $X_{S'}/S'$ 证明它们（对 (b) 用\Term{descent}{下降}；对 (iii) 用\Term{noetherian-induction}{诺特归纳}法及 3.4 (i)）。

最后考虑\Term{canonical}{典范}\Term{surjective}{满射}\Term{morphism}{态射} $\coprod X'_i\to X_{S'}$。由 4.3、3.5 和 3.7，$X_{S'}/S'$ 的 (b)、(iii) 可由各 $X'_i/S'$ 的 (c)、(iii) 得到。

这样，问题归结为在 $S$ \Term{integral}{整}、\Term{generic-point}{泛点}为 $s$，且 $X$ 在 $S$ 上\Term{projective}{射影}且\Term{integral}{整}的假设下证明 (c)、(iii)。取 $X/S$ 的\Term{very-ample}{甚丰沛}\Term{sheaf}{层} $\mathscr O_X(1)$。则对所有 $m,n$，$\chi(L_t^{\otimes m}(n))$ 及 $\chi(\mathscr O_{X_t}(n))$ 均与 $t\in S$ 无关（EGA III 7.9.4）。故 (c) 由 4.6 (i) $\Longleftrightarrow$ (v) 得到；这一\Term{equivalence}{等价}性还说明 $X/S$ 各\Term{fiber}{纤维}上 $\tau$-\Term{tau-equivalence}{等价}于 $0$ 的\Term{invertible-sheaf}{可逆层}的类所成的族只有一个
\SourcePage{660}
\Term{hilbert-polynomial}{Hilbert 多项式}，故 (iii) 由 2.11 得到。
\end{proof}
\textbf{注记 4.8。}设 $X$ 在\Term{noetherian}{诺特}\Term{scheme}{概形} $S$ 上\Term{proper}{固有}且\Term{integral}{整}。若不假设 $X/S$ \Term{projective}{射影}，尚不知道 $\operatorname{Pic}^\tau_{X/S}\to\operatorname{Pic}_{X/S}$ 是否仍为\Term{closed-immersion}{闭浸入}；但 $X$ 在 $S$ 上\Term{normal}{正规}且 $\operatorname{Pic}_{X/S}$ \Term{representable}{可表示}时除外（FGA 236，2.3）。

\SGAsection{5}{Néron--Severi 型有限性定理}
\begin{SGAresult}{\Term{theorem}{定理}}{5.1}
设 $X$ \Term{proper}{固有}于\Term{noetherian}{诺特}\Term{scheme}{概形} $S$。则 $X/S$ 的各\Term{geometric-fiber}{几何纤维} $X_K/K$ 的 \Term{neron-severi-group}{Néron--Severi 群}
\[\operatorname{Pic}_{X_K/K}(K)/\operatorname{Pic}^0_{X_K/K}(K)\]
均为\Term{finitely-generated}{有限生成}的，其秩及其\Term{torsion-subgroup}{挠子群} $\operatorname{Pic}^\tau_{X_K/K}(K)/\operatorname{Pic}^0_{X_K/K}(K)$ 的阶有统一的上界。
\end{SGAresult}
\begin{proof}
由\Term{noetherian-induction}{诺特归纳}法，只须将 $S$ 换为 $S_{\mathrm{red}}$ 的任意一个\Term{nonempty}{非空}\Term{open-subset}{开集}后证明 5.1；因而可设 $\operatorname{Pic}^\tau_{X/S}$ 可由\Term{finite-type}{有限型} $S$-\Term{scheme}{概形}表示（第 XII 讲 1.2 及 4.7 (iii)）。此时 $\operatorname{Pic}^\tau_{X_K/K}(K)/\operatorname{Pic}^0_{X_K/K}(K)$ 的阶等于 $(\operatorname{Pic}^\tau_{X/S})_K$ 的\Term{connected-component}{连通分支}数；在 $S$ 的某个\Term{nonempty}{非空}\Term{open-subset}{开集}之上，这个数保持\Term{invariant}{不变}（EGA IV 9.7.9）。

剩下只须给 $\operatorname{Pic}_{X_K/K}(K)/\operatorname{Pic}^0_{X_K/K}(K)$ 的秩一个上界。如同 4.7 (iii) 的证明，容易归结到 $S$ \Term{integral}{整}、$X$ 在 $S$ 上\Term{projective}{射影}且\Term{integral}{整}的情形。设 $r$ 为 $X/S$ 各\Term{fiber}{纤维}的共同\Term{dimension}{维数}。若 $r=0$（分别为 $1$），则这个秩显然为 $0$（分别为 $1$）。

若 $r\geq2$，我们将归结到 $r=2$ 且 $X$ 在 $S$ 上\Term{smooth}{光滑}的情形。缩小 $S$ 后，可找到 $X$ 的一个\Term{linear-section}{线性截面} $Y$，它在 $S$ 上仍\Term{integral}{整}，
\SourcePage{661}
各\Term{fiber}{纤维}的\Term{dimension}{维数}为 $2$；由 4.3，可将 $X$ 换为 $Y$。最后，由 4.3 及 3.10（使用\Term{surface}{曲面}的\Term{resolution-singularities}{奇点消解}），将 $S$ 换为 $S$ 的某个\Term{nonempty}{非空}\Term{open-subset}{开集}的\Term{flat}{平坦}\Term{covering}{覆盖}后，可将 $X$ 换为\Term{smooth}{光滑} $S$-\Term{scheme}{概形}。

现在，由 4.6 (i) $\Longleftrightarrow$ (iii) 以及\Term{etale-cohomology}{平展上同调}群的\Term{finiteness}{有限性}及\Term{specialization}{特化}\Term{theorem}{定理}（SGA 4，第 XVI 讲 2.2），5.1 可由下面的命题得到。
\end{proof}
\begin{SGAresult}{命题}{5.2}
设 $k$ 为\Term{algebraically-closed-field}{代数闭域}，$X$ 为\Term{proper}{固有}\Term{smooth}{光滑} $k$-\Term{scheme}{概形}。则 $X$ 上\Term{codimension}{余维}为 $p$ 的\Term{algebraic-cycle}{代数圈}模\Term{numerical-equivalence}{数值等价}所得的群 $C^p_{\mathrm{num}}(X)$ 是\Term{free-group}{自由群}，其秩不超过 Betti 数 $b^{2p}$。
\end{SGAresult}
\begin{proof}
若忘掉\Term{root-of-unity}{单位根}的\Term{action}{作用}，\Term{etale-cohomology}{平展上同调}
\[H^*(X,\mathbf Q_\ell)\mathrel{\mathop:}=\left[\varprojlim_\nu H^*(X,\mathbf Z/\ell^\nu\mathbf Z)\right]\otimes_{\mathbf Z_\ell}\mathbf Q_\ell,
\qquad (\ell,\operatorname{car}(k))=1,\]
便成为 Weil \Term{cohomology}{上同调}：它满足 \Term{poincare-duality}{Poincaré 对偶}及 \Term{kunneth-formula}{Künneth 公式}，并且接收一个来自\Term{algebraic-cycle}{代数圈}模\Term{rational-equivalence}{有理等价}所成的环 $C^*_{\mathrm{rat}}(X)$ 的非零\Term{functorial}{函子性}\Term{homomorphism}{同态}（见第 X 讲 7.13）。

设 $r=\dim X$，取 $a_1,\ldots,a_m\in C^{r-p}_{\mathrm{rat}}(X)$，使它们的像构成 $H^{2(r-p)}(X,\mathbf Q_\ell)$ 中由\Term{algebraic-cycle}{代数圈}生成的 $\mathbf Q_\ell$ \Term{vector-subspace}{向量子空间}的一组基。考虑\Term{homomorphism}{同态}
\[\alpha:C^p_{\mathrm{rat}}(X)\longrightarrow\mathbf Z^m,\qquad
\alpha(b)=(\langle b.a_1\rangle,\ldots,\langle b.a_m\rangle),\]
其中 $\langle\ \rangle$ 表示\Term{intersection-number}{交数}。容易看出 $\alpha$ 的像就是 $C^p_{\mathrm{num}}(X)$。
\end{proof}
\SourcePage{662}
\textbf{注记 5.3。}设 $k$ 为\Term{algebraically-closed-field}{代数闭域}，$X$ 为 $k$ 上的\Term{proper}{固有}\Term{smooth}{光滑}\Term{surface}{曲面}。第二 Betti 数 $b^2(X)$ \Term{solvable}{可解}释为 $c_2(X)-2+2b^1(X)$，其中 $b^1(X)=2\dim\operatorname{Pic}_{X/k}$，$c_2(X)$ 是第二\Term{chern-class}{Chern 类}，等于 \Term{euler-poincare-characteristic}{Euler--Poincaré 特征数}，也等于 $X\times X$ 的\Term{diagonal}{对角线}与自身的\Term{intersection-number}{交数}；因而 $b^2(X)$ 有一种不依赖\Term{cohomology}{上同调}的解释。在 $\operatorname{car}(k)>0$ 时，IGUSA 在\Term{etale-cohomology}{平展上同调}出现之前已证明了这种\Term{form}{形式}的不等式 $\operatorname{rang}(C^1_{\mathrm{num}}(X))\leq b^2(X)$；他使用了关于\Term{algebraic-curve}{代数曲线}删去\Term{finite}{有限}个点后的驯\Term{fundamental-group}{基本群}结构的深刻结果。

\begin{SGAresult}{\Term{theorem}{定理}}{5.4}
设 $k$ 为\Term{algebraically-closed-field}{代数闭域}，$X$ 为\Term{proper}{固有} $k$-\Term{scheme}{概形}。则存在 $X$ 中\Term{finite}{有限}条\Term{integral}{整}闭\Term{curve}{曲线} $Y_1,\ldots,Y_p$，具有如下性质：\Term{scheme}{概形} $\operatorname{Pic}_{X/k}$ 的一个\Term{subset}{子集} $\Lambda$ \Term{quasi-compact}{拟紧}，当且仅当 $\lambda\in\Lambda$ 的\Term{intersection-number}{交数} $\langle\lambda.Y_i\rangle$ \Term{bounded}{有界}。此外，可取 $p$ 为\Term{neron-severi-group}{Néron--Severi 群} $\operatorname{Pic}_{X/k}(k)/\operatorname{Pic}^0_{X/k}(k)$ 的秩。
\end{SGAresult}
\begin{proof}
由于 $\operatorname{Pic}^\tau_{X/k}$ 是 $\operatorname{Pic}_{X/k}$ 的\Term{quasi-compact}{拟紧}\Term{open-subgroup}{开子群}（4.7），$\Lambda$ \Term{quasi-compact}{拟紧}当且仅当\Term{vector-space}{向量空间}
\[[\operatorname{Pic}_{X/k}(k)/\operatorname{Pic}^\tau_{X/k}(k)]\otimes\mathbf Q\]
中对应的点集\Term{finite}{有限}。但 $\tau$-\Term{tau-equivalence}{等价}与\Term{numerical-equivalence}{数值等价}相同（4.6 (i) $\Longleftrightarrow$ (iii)），故可在 $X$ 的\Term{curve}{曲线}中选出若干条，使其像构成\Term{dual-space}{对偶空间}的一组基；断言由此得到。
\end{proof}
\SourcePage{663}
\SGAsection{6}{附录：研究 $P=\mathbf P^N_k$ 上的 (b)-层}
\textbf{6.1。}设 $k$ 为\Term{algebraically-closed-field}{代数闭域}，$P=\mathbf P^N_k$。设 $F$ 为 $P$ 上的\Term{coherent}{凝聚}\Term{sheaf}{层}，其\Term{hilbert-polynomial}{Hilbert 多项式}为
\[\chi(F(n))=\sum_{i=0}^r a_i\binom{n+i}{i},\]
并设 $(c)=(c_0,\ldots,c_r)$ 为逐项不小于 $(a)$ 的整数序列。令
\[a_{m,i}=\sum_{j=0}^{r-i}a_{j+i}\binom{m-1+j}{j},\qquad
c_{m,i}=\sum_{j=0}^{r-i}c_{j+i}\binom{m-1+j}{j}.\]
最后，对每个整数 $q\geq1$，以 $N_q$ 表示 $F$ 中\Term{support}{支集}\Term{dimension}{维数}不超过 $q-1$ 的最大\Term{coherent}{凝聚}子\Term{sheaf}{层}，并令 $F_q=F/N_q$。
\begin{SGAresult}{\Term{lemma}{引理}}{6.2}
在 6.1 的条件下，考虑\Term{exact-sequence}{正合列}
\[0\longrightarrow I\longrightarrow\mathscr O_P^{\oplus M}\longrightarrow F(m)\longrightarrow0,\qquad m\geq0.\]
则 $I$ 为 $p$-\Term{regular-sheaf}{正则层}，其中 $p=P_r(c_{m,0},\ldots,c_{m,r})$，$P_r$ 是第 $r$ 个 (b)-\Term{polynomial}{多项式}（1.9）。
\end{SGAresult}
\begin{proof}
显然 $\mathscr O_P$ 是 $(0,\ldots,0,1)$-\Term{b-sheaf}{层}（$N+1$ 项），故由 1.6，$I$ 是 $(0,\ldots,0,M)$-\Term{b-sheaf}{层}。由 2.10，
\[\chi(I(n))=M\binom{n+N}{N}-\sum_{i=0}^r a_{m,i}\binom{n+i}{i}.\]
因此断言由 1.11 及 1.10 得到。
\end{proof}
\begin{SGAresult}{命题}{6.3}
在 6.1 的条件下，设 $F$ 为 (b)-\Term{b-sheaf}{层}。则可用关于 $a_i,b_i$ 的通用\Term{polynomial}{多项式}，计算出整数 $m,M,m_1,M_1$，使有\Term{exact-sequence}{正合列}
\[\mathscr O_P(-m_1)^{\oplus M_1}\longrightarrow\mathscr O_P(-m)^{\oplus M}\longrightarrow F\longrightarrow0.\]
此外，计算 $M_1$ 时不涉及 $N$。
\end{SGAresult}
\SourcePage{664}
\begin{proof}
取 $m=P_{r-1}(c_0,\ldots,c_{r-1})$，其中 $c_i=b_i-a_i$，并取
\[M=\sum_i a_i\binom{m+i}{i},\]
则由 1.11 及 1.3，有\Term{exact-sequence}{正合列}
\[0\longrightarrow I\longrightarrow\mathscr O_P^{\oplus M}\longrightarrow F(m)\longrightarrow0.\]
同样，结合 6.2，取 $m_1=m+p$，其中 $p=P_r(a_{m,0},\ldots,a_{m,r})$，并取
\[M_1=M\binom{p+N}{N}-\sum_i a_i\binom{m_1+i}{i}=\chi(I(p)).\]
\end{proof}
\begin{SGAresult}{\Term{theorem}{定理}}{6.4}
在 6.1 的条件下，设 $F$ 为 $\mathscr O_P(-m)^{\oplus M}$ 的商，$m\geq0$。令
\[b_i=P_{r-i}(c_{m,i},\ldots,c_{m,r})\quad(i=0,\ldots,r),\]
其中 $P_s$ 为第 $s$ 个 (b)-\Term{polynomial}{多项式}（1.9）；令 $b=b_1+m-1$，并令
\[B=\sum_{i=1}^r a_i\binom{b+i}{i}.\]
则：
\begin{enumerate}
\item[(i)] $F$ 为 $b$-\Term{regular-sheaf}{正则层}；
\item[(ii)] $-B\leq a_0=h^0(F(b))-B$；
\item[(iii)] $F$ 为 $(b_0,\ldots,b_r)$-\Term{b-sheaf}{层}。
\end{enumerate}
\end{SGAresult}
\begin{proof}
取 $\mathscr O_P(1)$ 的\Term{section}{截面} $\sigma$，使得到如下行列均\Term{exact-sequence}{正合}的\Term{commutative-diagram}{交换图}：
\[
\begin{tikzcd}[column sep=large,row sep=small]
&0\arrow[d]&0\arrow[d]&0\arrow[d]&\\
0\arrow[r]&I(-1)\arrow[r]\arrow[d]&I\arrow[r]\arrow[d]&J\arrow[r]\arrow[d]&0\\
0\arrow[r]&\mathscr O_P(-1)^{\oplus M}\arrow[r]\arrow[d]&\mathscr O_P^{\oplus M}\arrow[r]\arrow[d]&\mathscr O_Y^{\oplus M}\arrow[r]\arrow[d]&0\\
0\arrow[r]&F(m-1)\arrow[r]\arrow[d]&F(m)\arrow[r]\arrow[d]&G(m)\arrow[r]\arrow[d]&0\\
&0&0&0&
\end{tikzcd}
\]
\SourcePage{665}
由 1.7，$\chi(G(n))=\sum_{i=0}^{r-1}a_{i+1}\binom{n+i}{i}$，故由 6.2，$J$ 为 $b_1$-\Term{regular-sheaf}{正则层}；因此，当 $q\geq2$ 且 $p\geq b_1$ 时，由 1.4 (i) 得 $H^q(I(p))=0$。\Term{exact-sequence}{正合列}
\[H^q(\mathscr O_P^{\oplus M}(p))\longrightarrow H^q(F(m+p))\longrightarrow H^{q+1}(I(p))\]
遂说明 $F(m)$ 为 $(b_1-1)$-\Term{regular-sheaf}{正则层}，这就证明了 (i)。

由 (i)，
\[0\leq h^0(F(b))=\chi(F(b))=\sum_{i=0}^r a_i\binom{b+i}{i},\]
从而得到 (ii)。由 2.10 及 1.9，
\[\chi(F(b))=\sum_{i=0}^r a_{m,i}\binom{b_1-1+i}{i}
\leq b_1+\sum_{i=0}^r c_{m,i}\binom{b_1-1+i}{i}=b_0;\]
故 $h^0(F(-1))\leq h^0(F(b))\leq b_0$。对 $r$ 归纳，可设 $G$ 是 $(b_1,\ldots,b_r)$-\Term{b-sheaf}{层}，从而由 1.6.2 得到 (iii)。
\end{proof}
\begin{SGAresult}{\Term{lemma}{引理}}{6.5}
在 6.1 的条件下，设 $F$ 没有\Term{dimension}{维数}为 $0$ 的\Term{associated-prime-cycle}{伴随素圈}。
\begin{enumerate}
\item[(i)] 若 $h^0(F)\geq1$，则 $h^0(F(-1))\leq h^0(F)-1$。
\item[(ii)] 当 $n\geq h^0(F)$ 时，$H^0(F(-n))=0$。
\item[(iii)] 设有\Term{exact-sequence}{正合列} $0\to F(-1)\to F\to G\to0$，且整数 $n_0>0$ 满足 $H^0(G(-n_0))=0$；则 $H^0(F(-n_0))=0$，且 $h^0(F)\leq n_0h^0(G)$。
\end{enumerate}
\end{SGAresult}
\begin{proof}
取 $\sigma\in H^0(F)$，$\sigma\neq0$；则 $\sigma$ 定义一个非零子\Term{sheaf}{层} $F'=\mathscr O_{X'}$，其\Term{support}{支集}\Term{dimension}{维数} $s=\dim\operatorname{supp}(F')\geq1$。取如下行列均\Term{exact-sequence}{正合}的\Term{commutative-diagram}{交换图}：
\[
\begin{tikzcd}[row sep=small]
&0\arrow[d]&0\arrow[d]&0\arrow[d]&\\
0\arrow[r]&F'(-1)\arrow[r]\arrow[d]&F'\arrow[r]\arrow[d]&G'\arrow[r]\arrow[d]&0\\
0\arrow[r]&F(-1)\arrow[r]&F\arrow[r]&G\arrow[r]&0
\end{tikzcd}
\]
\SourcePage{666}
由 $s\geq1$ 得 $G'\neq0$，故 $F'$ 不包含于 $F(-1)$；也就是说，\Term{section}{截面} $\sigma$ 不属于\Term{injectivity}{单射} $H^0(F(-1))\hookrightarrow H^0(F)$ 的像，故得 (i)。

(ii) 立即由 (i) 得到，只须注意 $h^0(F)=0$ 显然蕴含 $h^0(F(-1))=0$，因为 $h^0(F(-1))\leq h^0(F)$。

在 (iii) 中，有 $0\leq h^0(F(n))-h^0(F(n-1))\leq h^0(G(n))$。因此当 $n\geq n_0$ 时，由 $h^0(G(-n_0))=0$ 得 $h^0(F(-n_0))=h^0(F(-n))$；再由 (ii) 得 $h^0(F(-n_0))=0$。从而，对 $p\geq1$，
\[h^0(F(p-n_0))\leq h^0(G(1-n_0))+\cdots+h^0(G(p-n_0))\leq p h^0(G(p-n_0)),\]
这就证明了 (iii)。
\end{proof}
\begin{SGAresult}{命题}{6.6}
在 6.1 的条件下，设 $F$ 为 $(b_0,\ldots,b_r)$-\Term{b-sheaf}{层}。则 $F_q$ 是具有如下参数的\Term{b-sheaf}{层}：
\[((b_{q-1})^q,\ldots,(b_{q-1})^2,b_{q-1},b_q,\ldots,b_r).\]
\end{SGAresult}
\begin{proof}
取 $\mathscr O_X(1)$ 的一个“\Term{generic}{一般}”\Term{section}{截面} $\sigma$；它定义如下行列均\Term{exact-sequence}{正合}的\Term{commutative-diagram}{交换图}：
\[
\begin{tikzcd}[row sep=small]
&0\arrow[d]&0\arrow[d]&0\arrow[d]&\\
0\arrow[r]&N_q(-1)\arrow[r]\arrow[d]&N_q\arrow[r]\arrow[d]&G'\arrow[r]\arrow[d]&0\\
0\arrow[r]&F(-1)\arrow[r]\arrow[d]&F\arrow[r]\arrow[d]&G\arrow[r]\arrow[d]&0\\
0\arrow[r]&F_q(-1)\arrow[r]\arrow[d]&F_q\arrow[r]\arrow[d]&G''\arrow[r]\arrow[d]&0\\
&0&0&0&
\end{tikzcd}
\]
其中 $G$ 为 $(b_1,\ldots,b_r)$-\Term{b-sheaf}{层}（1.6 (i)），且 $G''\simeq G_{q-1}$（由\Term{depth}{深度}的考虑）。
\SourcePage{667}
先设 $q=1$。此时 $\dim\operatorname{supp}(N_q)=0$，故 $H^1(N_q(-1))=0$，且 $h^0(F_q(-1))\leq h^0(F(-1))\leq b_0$。又 $G'=0$，故 $G\simeq G''$。因此 $F_1$ 确为 $(b_0,\ldots,b_r)$-\Term{b-sheaf}{层}。此外，6.5 (ii) 说明 $H^0(F_1(-b_0))=0$。若 $q\geq2$，对 $q$ 归纳，可设 $G_{q-1}$ 是 $((b_{q-1})^{q-1},\ldots,b_{q-1},\ldots,b_r)$-\Term{b-sheaf}{层}，且 $H^0(G_{q-1}(-b_{q-1}))=0$。于是 6.5 说明 $H^0(F_q(-b_{q-1}))=0$ 且 $h^0(F_q(-1))\leq(b_{q-1})^q$，断言由此得到。
\end{proof}
\begin{SGAresult}{\Term{theorem}{定理}}{6.7}
存在两族\Term{polynomial}{多项式} $\{A_i(X_0,\ldots,X_i;Y)\}$ 及 $\{A_i^{(q)}(X_0,\ldots,X_q;Y)\}$，具有如下性质。假设满足 6.1 的条件，且 $F$ 为 $\mathscr O_P(-m)^{\oplus M}$ 的商，其中 $m\geq0$。
\begin{enumerate}
\item[(i)] 若 $F$ 是 $(b_0,\ldots,b_r)$-\Term{b-sheaf}{层}，则对 $i=0,\ldots,r$，
\[|a_i|\leq A_{r-i}(b_i,\ldots,b_r;m).\]
\item[(ii)] 若 $\chi(F_q(n))=\sum_{i=0}^r a_i^{(q)}\binom{n+i}{i}$ 为 $F_q$ 的\Term{hilbert-polynomial}{Hilbert 多项式}，则对 $i=0,\ldots,q-1$，
\[|a_i^{(q)}|\leq A_{r-i}^{(r-q)}(c_{q-1},\ldots,c_r;m).\]
\end{enumerate}
（注意 $a_{q-1}^{(q)}\leq a_{q-1}$，且 $a_q^{(q)}=a_q,\ldots,a_r^{(q)}=a_r$，因为 $\dim\operatorname{supp}(N_q)\leq q-1$。）
\end{SGAresult}
\begin{proof}
取\Term{exact-sequence}{正合列} $0\to F(-1)\to F\to G\to0$，使 $G$ 为 $(b_1,\ldots,b_r)$-\Term{b-sheaf}{层}。对 $r$ 归纳，结合 1.7，可设 $A_0,\ldots,A_{r-1}$ 已定义。此时，6.4 (ii) 结合 1.8 即给出 $A_r$ 的定义。最后，(ii) \Term{form}{形式}地由 6.4 (iii)、6.6 及 (i)（在其中以 $F_q$ 代 $F$）得到。
\end{proof}
\SourcePage{668}
\begin{SGAresult}{推论}{6.8（Grothendieck）}
设 $X$ \Term{projective}{射影}于\Term{noetherian}{诺特}\Term{scheme}{概形} $S$，$\mathscr O_X(1)$ 对 $X/S$ \Term{very-ample}{甚丰沛}，$\mathscr F$ 为 $X/S$ 各\Term{fiber}{纤维}上的\Term{coherent}{凝聚}\Term{sheaf}{层}的类所成的族。假设：
\begin{enumerate}
\item[(a)] 存在 $X$ 上的\Term{coherent}{凝聚}\Term{sheaf}{层} $E$，使 $\mathscr F$ 包含在形如 $E_K$ 的\Term{sheaf}{层}的商的类所成的族中；
\item[($\mathrm b_q$)] 各 $F_K\in\mathscr F$ 的\Term{hilbert-polynomial}{Hilbert 多项式} $\chi(F_K(n))$ 中，\Term{degree}{次数}不小于 $q-1$ 的各项系数有统一的上界。
\end{enumerate}
则当 $F_K$ 遍历 $\mathscr F$ 时，各 $(F_K)_q$ 构成\Term{bounded-family}{有界族}；此外，$\chi(F_K(n))$ 中\Term{degree}{次数}不小于 $q-2$ 的各项系数有统一的下界。
\end{SGAresult}
\begin{proof}
显然可设 $X=\mathbf P^N_S$ 且 $E=\mathscr O_X(-m)^{\oplus M}$。此时，第一项断言\Term{form}{形式}地由 6.7 (ii) 及 1.13 得到；第二项则用\Term{exact-sequence}{正合列} $0\to F(-1)\to F\to G\to0$ 归纳，并由 6.4 (ii) 得到。
\end{proof}
\begin{SGAresult}{定义}{6.9}
设 $k$ 为域。称配有\Term{very-ample}{甚丰沛}\Term{sheaf}{层} $\mathscr O_X(1)$ 的\Term{projective}{射影} $k$-\Term{scheme}{概形} $X$ 为\Term{dimension}{维数} $r$ 的\Term{special-positive-cycle}{特殊正 $k$-圈}（$k$-cycle positif spécial），若 $X$ 是若干\Term{dimension}{维数} $r$ 的闭子\Term{scheme}{概形} $X_j$ 的并，且每个 $X_j$ 都是某个\Term{integral}{整} $k_j$-\Term{scheme}{概形} $X'_j$（配有 $\mathscr O_{X'_j}(1)$）经\Term{base-change}{基变换} $\operatorname{Spec}(k)\to\operatorname{Spec}(k_j)$ 得到的\Term{inverse-image}{逆像}。称\Term{hilbert-polynomial}{Hilbert 多项式}
\[\chi(\mathscr O_X(n))=\sum a_i\binom{n+i}{i}\]
中的系数 $a_r$ 为 $X$ 的\Term{degree}{次数}。
\end{SGAresult}
\begin{SGAresult}{\Term{lemma}{引理}}{6.10}
设 $k$ 为\Term{algebraically-closed-field}{代数闭域}，$X$ 为\Term{dimension}{维数} $r$、\Term{degree}{次数} $d$ 的\Term{special-positive-cycle}{特殊正 $k$-圈}。则 $\mathscr O_X$ 是 $(0,\ldots,0,d)$-\Term{b-sheaf}{层}。
\end{SGAresult}
\SourcePage{669}
\begin{proof}
设 $X\subset\mathbf P^N_k$。由 1.6 (i)，可将 $k$ 换为一个\Term{purely-transcendental-extension}{纯超越扩张}的\Term{algebraic-closure}{代数闭包}；因而可设 $k$ 中存在 $r(N+1)$ 个在所有 $k_j$ 上\Term{algebraic-independence}{代数独立}的元素。则由这些元素定义的 $X$ 的\Term{continuous}{连续} $r$ 个“泛”\Term{hyperplane-section}{超平面截面}，仍然都是\Term{degree}{次数} $d$ 的\Term{special-positive-cycle}{特殊正 $k$-圈}。

对 $r$ 归纳，只须注意：当 $r=0$ 时，$h^0(\mathscr O_X(-1))=d$；当 $r\geq1$ 时，$h^0(\mathscr O_X(-1))=0$。对于后一结论，由\Term{injectivity}{单射} $\mathscr O_X\hookrightarrow\prod_j\mathscr O_{X_j}$，可设 $X$ 本身来自某个\Term{integral}{整}\Term{scheme}{概形} $X'$；但此时 $h^0(\mathscr O_X(-1))=h^0(\mathscr O_{X'}(-1))$，后者为 $0$，因为 $H^0(\mathscr O_{X'})$ 是域。
\end{proof}
\begin{SGAresult}{推论}{6.11}
\begin{enumerate}
\item[(i)] 设 $r,d$ 为两个整数，对 $i=0,\ldots,r$ 令 $c_i=A_{r-i}(0,\ldots,0,d;0)$（6.7 (i)），并令 $p=P_r(c_0,\ldots,c_r)$（1.9）。设 $k$ 为\Term{algebraically-closed-field}{代数闭域}，$X$ 为\Term{dimension}{维数}不超过 $r$、\Term{degree}{次数}不超过 $d$ 的\Term{special-positive-cycle}{特殊正 $k$-圈}，其\Term{hilbert-polynomial}{Hilbert 多项式}为 $\chi(\mathscr O_X(n))=\sum_{i=0}^r a_i\binom{n+i}{i}$。则对 $i=0,\ldots,r$，
\[|a_i|\leq c_i.\]
此外，$X$ 可\Term{immersion}{浸入} $\mathbf P^N_k$ 为子\Term{scheme}{概形}，其中 $N=d(r+1)-1$，并由至多 $\binom{N+p}{N}$ 个 $p$ 次方程定义。
\item[(ii)] （Chow。）设 $S$ 为\Term{noetherian}{诺特}\Term{scheme}{概形}。当 $K$ 遍历各 $k(s)$（$s\in S$）的\Term{algebraically-closed}{代数闭}\Term{group-extension}{扩张}时，\Term{dimension}{维数}及\Term{degree}{次数}\Term{bounded}{有界}的\Term{special-positive-cycle}{特殊正 $K$-圈}构成一个“\Term{bounded-family}{有界族}”（可将 $X$ 看作某个固定的 $\mathbf P^N_S$ 的子\Term{scheme}{概形}，也可取其显然的抽象意义）。
\end{enumerate}
\end{SGAresult}
\SourcePage{670}
\begin{proof}
在 (i) 中，第一项断言由 6.7 (i) 得到；第二项由 6.10、1.8 (i)，以及应用于 $0\to I\to\mathscr O_P\to\mathscr O_X\to0$（$P=\mathbf P^N_k$）的 6.2 和 1.3 得到。最后，(ii) 由 (i) 及 1.13 得到。
\end{proof}
\textbf{注记 6.12。}设 $k$ 为\Term{algebraically-closed-field}{代数闭域}，$X$ 为纯 $r$ 维的\Term{projective}{射影} $k$-\Term{scheme}{概形}，没有\Term{embedding}{嵌入}的\Term{associated-prime-cycle}{伴随素圈}。若容许 $\mathscr O_X$ 中有任意的\Term{nilpotent}{幂零}元素，便不能只用 $a_r$ 来界定\Term{hilbert-polynomial}{Hilbert 多项式} $\chi(\mathscr O_X(n))=\sum_{i=0}^r a_i\binom{n+i}{i}$ 中的各系数。

例如，设 $Z$ 为 $\mathbf P^3_k$ 中一条\Term{degree}{次数} $d$ 的\Term{smooth}{光滑}\Term{curve}{曲线}。对每个 $n\gg0$，存在一个包含 $Z$ 的 $n$ 次\Term{smooth}{光滑}\Term{surface}{曲面} $Y$。于是 $Z$ 是 $Y$ 上的\Term{cartier-divisor}{Cartier 除子}；令 $\mathscr O_{X_n}=\mathscr O_Y/\mathscr O_Y(-2Z)$。最后，有\Term{exact-sequence}{正合列}
\[0\longrightarrow\mathscr O_X(-Z)\longrightarrow\mathscr O_X\longrightarrow\mathscr O_Z\longrightarrow0\]
及公式 $\langle Z^2\rangle_Y=2p_a(Z)-2-(n-4)d$，从而
\[a_0=\chi(\mathscr O_{X_n}(-1))\longrightarrow\infty\qquad(n\to\infty),\]
而 $a_1=2d$。
\begin{SGAresult}{推论}{6.13}
在 6.1 的条件下，考虑\Term{exact-sequence}{正合列}
\[\mathscr O_P(-m_1)^{\oplus M_1}\longrightarrow\mathscr O_P(-m)^{\oplus M}\xrightarrow{\alpha}F\longrightarrow0,\]
其中 $m_1\geq m\geq0$。则对 $i=0,\ldots,N$，
\[|a_i|\leq A_{N-i}(0,\ldots,0,M;m_1)+M\binom{m}{N-i}.\]
\end{SGAresult}
\begin{proof}
令 $I=\operatorname{Ker}(\alpha)$。显然 $I$ 为 $(0,\ldots,0,M)$-\Term{b-sheaf}{层}，且
\[\chi(I(n))=\sum_{i=0}^N\left[M\binom{-m-1+N-i}{N-i}-a_i\right]\binom{n+i}{i}.\]
故断言由 6.7 (i) 及公式
\[\binom{-m-1+Q}{Q}=(-1)^Q\binom{m}{Q}\]
得到。
\end{proof}
\SourcePage{671}
\begin{SGAresult}{推论}{6.14（I. Hermann）}
对每个 $N\geq0$，存在\Term{polynomial}{多项式} $R_N(m)$，使对任意域 $k_0$ 及\Term{polynomial-ring}{多项式环} $k_0[T_1,\ldots,T_N]$ 中由\Term{degree}{次数}不超过 $m$ 的元素生成的任意\Term{ideal}{理想} $I$，其\Term{radical-ideal}{根理想} $\sqrt I$ 均由\Term{degree}{次数}不超过 $R_N(m)$ 的元素生成。
\end{SGAresult}
\begin{proof}
取 $k$ 为 $k_0$ 的\Term{algebraic-closure}{代数闭包}。引入一个辅助\Term{indeterminate}{不定元}，令 $Y$（分别为 $X$）为 $P=\mathbf P^N_k$ 中由 $I$（分别为 $\sqrt I$）的\Term{homogeneous}{齐次}化\Term{ideal}{理想}定义的子\Term{scheme}{概形}。在\Term{hilbert-polynomial}{Hilbert 多项式}
\[\chi(\mathscr O_Y(n))=\sum_{i=0}^N b_i\binom{n+i}{i}\]
中，由 6.13 有 $|b_i|\leq c_i$，其中 $c_i=A_{N-i}(0,\ldots,0,1;m)$。

设 $X=\bigcup X^q$ 为按\Term{dimension}{维数}的\Term{resolution}{分解}。则 $\deg(X^q)\leq e_q$，其中 $e_q=P_{N-q}(c_q,\ldots,c_N)$。事实上，以一个\Term{codimension}{余维}为 $q$ 的“\Term{generic}{一般}”\Term{linear}{线性}\Term{space}{空间}截 $X$（分别为 $Y$）后，可设 $q=0$。此时显然 $\deg(X^q)=h^0(\mathscr O_{X^q})$ 且 $h^0(\mathscr O_{X^q})\leq h^0(\mathscr O_Y)$，而由 6.4 (iii)，$h^0(\mathscr O_Y)\leq e_q$。

因此，由 6.10，$\mathscr O_{X^q}$ 为 $(0,\ldots,0,e_q)$-\Term{b-sheaf}{层}。由\Term{injectivity}{单射} $\mathscr O_X\hookrightarrow\prod_q\mathscr O_{X^q}$ 可知 $\mathscr O_X$ 为 $(e_0,\ldots,e_N)$-\Term{b-sheaf}{层}。因而，由 6.7，在\Term{hilbert-polynomial}{Hilbert 多项式} $\chi(\mathscr O_X(n))=\sum_{i=0}^n a_i\binom{n+i}{i}$ 中，有 $|a_i|\leq f_i$，其中 $f_i=A_{N-i}(e_i,\ldots,e_N;0)$。最后，$X$ 的\Term{ideal}{理想} $\sqrt I$ 为 $R_N(m)$-\Term{regular-sheaf}{正则层}，其中 $R_N(m)=P_N(f_0,\ldots,f_N)$；由 1.3 即得断言。
\end{proof}
\SourcePage{672}
\SGAsection{7}{附录：Hodge 指数定理}
这里用直接的论证重新得到并补充前面的一些结果。
\begin{SGAresult}{\Term{theorem}{定理}}{7.1}
设 $X$ 为\Term{algebraically-closed-field}{代数闭域}上的\Term{proper}{固有}\Term{irreducible}{不可约}\Term{surface}{曲面}。假设 $X$ 上存在\Term{invertible-sheaf}{可逆层} $H$，使 $\langle c_1(H)^2\rangle>0$（例如 $H$ \Term{ample}{丰沛}）。则 $X$ 上满足
\[\langle c_1(L).c_1(H)\rangle=0,\qquad\langle c_1(L)^2\rangle\geq0\]
的\Term{invertible-sheaf}{可逆层} $L$ \Term{numerical-equivalence}{数值等价}于 $0$。（因此，$A^1_{\mathrm{num}}(X)_{\mathbf Q}$ 上的\Term{intersection-form}{交形式}为\Term{nondegenerate}{非退化}的，型为 $(1,\rho-1)$；由 4.6，$A^1_{\mathrm{num}}(X)_{\mathbf Q}=[\operatorname{Pic}(X)/\operatorname{Pic}^\tau(X)]\otimes\mathbf Q$。）
\end{SGAresult}
\begin{proof}
\textbf{7.1.1。}先假设：$X$ \Term{reduced}{约化}（因而\Term{integral}{整}），$H$ \Term{ample}{丰沛}，$\langle c_1(L).c_1(H)\rangle=0$，且 $\langle c_1(L)^2\rangle>0$。用反证法，取 $m>0$ 使 $H_1=L\otimes H^{\otimes m}$ \Term{ample}{丰沛}；再取 $n$，使令 $N=L^{\otimes n}\otimes H^{\otimes-1}$ 时，
\[\langle c_1(N).c_1(H_1)\rangle=n\langle c_1(L)^2\rangle-m\langle c_1(H)^2\rangle>0.\]
此时还有 $\langle c_1(N).c_1(H)\rangle=-\langle c_1(H)^2\rangle<0$，以及 $\langle c_1(N)^2\rangle=n^2\langle c_1(L)^2\rangle+\langle c_1(H)^2\rangle>0$，这与下面的\Term{lemma}{引理}矛盾。
\end{proof}
\begin{SGAresult}{\Term{lemma}{引理}}{7.1.2}
设 $X$ 为\Term{algebraically-closed-field}{代数闭域}上的\Term{integral}{整}\Term{projective}{射影}\Term{surface}{曲面}。对 $X$ 上满足 $\langle c_1(N)^2\rangle>0$ 的\Term{invertible-sheaf}{可逆层} $N$，以下条件\Term{equivalence}{等价}：
\begin{enumerate}
\item[(i)] 对每个\Term{ample}{丰沛}\Term{sheaf}{层} $H$，$\langle c_1(N).c_1(H)\rangle>0$；
\item[(i')] 存在\Term{ample}{丰沛}\Term{sheaf}{层} $H$，使 $\langle c_1(N).c_1(H)\rangle>0$；
\item[(ii)] 存在整数 $n>0$，使 $H^0(X,N^{\otimes n})\neq0$。
\end{enumerate}
\end{SGAresult}
\SourcePage{673}
\begin{proof}
(ii) 蕴含存在\Term{effective-divisor}{有效除子} $D$ 使 $N^{\otimes n}\simeq\mathscr O_X(D)$，而 $\langle c_1(N)^2\rangle>0$ 蕴含 $D>0$，故得 (i)。

为证明 (i') 蕴含 (ii)，显然可设 $H=\mathscr O_X(1)$ \Term{very-ample}{甚丰沛}。于是存在\Term{integral}{整}、\Term{locally-principal}{局部主}闭\Term{curve}{曲线} $Y$，使 $H=\mathscr O_X(Y)$。由 Riemann \Term{theorem}{定理}，(i') 蕴含当 $p\geq0$ 且 $n\geq n_0$ 时，$H^1(N_Y^{\otimes n}(p))=0$，其中 $n_0=2h^1(\mathscr O_Y)-1$。由于 $p\gg0$ 时 $H^2(N^{\otimes n}(p))=0$，\Term{exact-sequence}{正合列}
\[H^1(N_Y^{\otimes n}(p))\longrightarrow H^2(N^{\otimes n}(p-1))\longrightarrow H^2(N^{\otimes n}(p))\longrightarrow0\]
说明当 $n\geq n_0$ 时 $H^2(N^{\otimes n})=0$。因此当 $n\gg0$ 时，
\[h^0(N^{\otimes n})=\chi(N^{\otimes n})+h^1(N^{\otimes n})
\geq\chi(N^{\otimes n})=\tfrac12\langle c_1(N)^2\rangle n^2+\cdots>0,\]
\Term{lemma}{引理}得证。
\end{proof}
\begin{proof}[\Term{theorem}{定理} 7.1 的证明（续）]
\textbf{7.1.3。}现在在 7.1 中假设存在 $X$ 上的\Term{invertible-sheaf}{可逆层} $M$，使 $\langle c_1(L).c_1(M)\rangle\neq0$，并用反证法。将 $X$ 换为 $X$ 的某个 \Term{chow-presentation}{Chow 表示}的底层\Term{reduced}{约化}\Term{scheme}{概形}，由\Term{projection-formula}{投影公式}（第 X 讲 4.4.4）知这是允许的。取 $X$ 上的\Term{ample}{丰沛}\Term{sheaf}{层} $G$。取 $m>0$，使 $H_1=G^{\otimes m}\otimes M$ \Term{ample}{丰沛}且 $\langle c_1(L).c_1(H_1)\rangle\neq0$；再取 $p,q$（$q\neq0$），使令 $L_1=L^{\otimes p}\otimes H^{\otimes q}$ 时，
\[\langle c_1(L_1).c_1(H_1)\rangle
=p\langle c_1(L).c_1(H_1)\rangle+q\langle c_1(H).c_1(H_1)\rangle=0.\]
此时还有
\[\langle c_1(L_1)^2\rangle=p^2\langle c_1(L)^2\rangle+q^2\langle c_1(H)^2\rangle>0,\]
从而归结为 7.1.1 的情形，其中以 $L_1$ 代 $L$，以 $H_1$ 代 $H$。

\textbf{7.1.4。}最后考虑 7.1 的\Term{generic}{一般}情形。仍用反证法，假设 $L$ 不\Term{numerical-equivalence}{数值等价}于 $0$；再将 $X$ 换为 $X_{\mathrm{red}}$ 的\Term{normalization}{正规化}，使 $X$ 只有
\SourcePage{674}
\Term{finite}{有限}个\Term{singular-point}{奇点}；由\Term{projection-formula}{投影公式}（第 X 讲 4.4.4）及 4.5 (ii) 可知这一归约是合理的。因此 $X$ 上存在一条\Term{integral}{整}闭\Term{curve}{曲线} $Y$，使 $\langle c_1(L).Y\rangle\neq0$，且 $Y$ 在某个\Term{finite}{有限}集 $Z$ 之外\Term{locally-principal}{局部主}。取 $\mathscr O_Y$ 的\Term{ideal}{理想} $I$ 的\Term{blow-up}{爆破} $f:X'\to X$。\Term{ideal}{理想} $M=I\mathscr O_{X'}$ \Term{invertible}{可逆}，定义 $X'$ 的一个子\Term{scheme}{概形}，它是两个子\Term{scheme}{概形} $Y',Y''$ 的并，满足 $f(Y')=Y$、$f(Y'')\subset Z$。由\Term{projection-formula}{投影公式}（第 X 讲 4.4.4），
\[\langle c_1(f^*L).c_1(M)\rangle
=-\langle c_1(f^*L).Y'\rangle-\langle c_1(f^*L).Y''\rangle
=-\langle c_1(L).Y\rangle\neq0.\]
可将 $X,L,H$ 换为 $X',f^*L,f^*H$；于是归结为前一情形，7.1 的证明完毕。
\end{proof}
\textbf{注记 7.2。}一个\Term{algebraically-closed-field}{代数闭域}上的\Term{proper}{固有}\Term{irreducible}{不可约}\Term{surface}{曲面}若具有满足 $\langle c_1(H)^2\rangle>0$ 的\Term{invertible-sheaf}{可逆层} $H$，是否必然\Term{projective}{射影}，尚属未知。对于每个\Term{proper}{固有}\Term{irreducible}{不可约}\Term{surface}{曲面}，$A^1_{\mathrm{num}}(X)_{\mathbf Q}$ 上的\Term{intersection-form}{交形式}是否\Term{nondegenerate}{非退化}，也尚属未知；由 7.1 可知，若它\Term{degeneration}{退化}，便必为负的（非定号）\Term{form}{形式}。
\begin{SGAresult}{推论}{7.3}
在 7.1 的条件下，设 $M$ 为 $X$ 上不\Term{numerical-equivalence}{数值等价}于 $H$ 的\Term{invertible-sheaf}{可逆层}，并令 $d=\langle c_1(M).c_1(H)\rangle$、$h=\langle c_1(H)^2\rangle$。则
\[\langle c_1(M)^2\rangle<d^2/h.\]
\end{SGAresult}
\begin{proof}
令 $L=M^{\otimes n}\otimes H^{\otimes-d}$。由于 $\langle c_1(L).c_1(H)\rangle=0$，由 7.1 得 $\langle c_1(L)^2\rangle<0$，断言由此得到。
\end{proof}
\SourcePage{675}
\begin{SGAresult}{推论}{7.4}
设 $X$ 为\Term{algebraically-closed-field}{代数闭域}上的\Term{irreducible}{不可约}\Term{projective}{射影}\Term{scheme}{概形}，\Term{dimension}{维数} $r\geq2$，$H$ 为 $X$ 上的\Term{ample}{丰沛}\Term{sheaf}{层}。则：
\begin{enumerate}
\item[(i)] $X$ 上的\Term{invertible-sheaf}{可逆层} $L$ \Term{numerical-equivalence}{数值等价}于 $0$ 的（必要且）充分条件是
\[\langle c_1(L).c_1(H)^{r-1}\rangle=0,\qquad
\langle c_1(L)^2.c_1(H)^{r-2}\rangle\geq0.\]
\item[(ii)] 对 $X$ 上每个不\Term{numerical-equivalence}{数值等价}于 $H$ 的\Term{invertible-sheaf}{可逆层} $M$，有
\[\langle c_1(M)^2.c_1(H)^{r-2}\rangle
<\frac{\langle c_1(M).c_1(H)^{r-1}\rangle^2}{\langle c_1(H)^r\rangle}.\]
\end{enumerate}
\end{SGAresult}
\begin{proof}
根据 7.3 的论证，(ii) 立即由 (i) 得到。为证明 (i)，取 $X$ 中的一条\Term{integral}{整}闭\Term{curve}{曲线} $Y$，须证明 $\langle c_1(L).Y\rangle=0$。用 $H^{\otimes n}$（$n>0$）将 $X$ \Term{embedding}{嵌入} $\mathbf P^N_k$。设 $I=\bigoplus I_\nu$ 为 $Y$ 的\Term{homogeneous-ideal}{齐次理想}，取整数 $\nu_0$，使对 $\nu\geq\nu_0$，$I_{\nu_0}$ 生成 $I_\nu$。于是，包含 $Y$ 的 $\nu_0$ 次\Term{hypersurface}{超曲面} $Z$ 在 $Y$ 之外没有公共点。对 $r=\dim X$ 归纳，借助 Bertini \Term{theorem}{定理}，可找到包含 $Y$ 的 $r-2$ 个 $\nu_0$ 次\Term{hypersurface}{超曲面} $Z_1,\ldots,Z_{r-2}$，使 $X'=Z_1\cap\cdots\cap Z_{r-2}\cap X$ \Term{irreducible}{不可约}。由于 $\langle c_1(L).Y\rangle=\langle c_1(L_{X'}).Y\rangle$，只须证明 $L_{X'}$ 在 $X'$ 上\Term{numerical-equivalence}{数值等价}于 $0$。而由\Term{projection-formula}{投影公式}，有 $\langle c_1(L_{X'}).c_1(H_{X'})\rangle=0$ 及 $\langle c_1(L_{X'})^2\rangle\geq0$，故断言由 7.1 得到。
\end{proof}
\begin{SGAresult}{推论}{7.5}
设 $k$ 为\Term{algebraically-closed-field}{代数闭域}，$Y$ 为配有\Term{ample}{丰沛}\Term{sheaf}{层} $H$ 的\Term{projective}{射影} $k$-\Term{scheme}{概形}，$X$ 为 $H$ 的一个\Term{section}{截面} $\varphi$ 的零点\Term{scheme}{概形}。假设对 $Y$ 的每个\Term{irreducible-component}{不可约分支} $Y^i$，
\SourcePage{676}
$X\cap Y^i=V(\varphi|_{Y^i})$ 的各\Term{irreducible-component}{不可约分支}的\Term{dimension}{维数}均不小于 $2$。则\Term{morphism}{态射} $A^1_{\mathrm{num}}(Y)\to A^1_{\mathrm{num}}(X)$ 是\Term{injectivity}{单射}。
\end{SGAresult}
\begin{proof}
显然可设 $Y$ \Term{integral}{整}；于是借助\Term{projection-formula}{投影公式}，断言由 7.4 (i) 得到。
\end{proof}

\textbf{参考文献}
\begin{enumerate}
\item[{[1]}] Abhyankar S.：《算术\Term{surface}{曲面}的\Term{resolution-singularities}{奇点消解}》（Resolution of singularities of arithmetical surfaces），载《算术\Term{algebraic-geometry}{代数几何}》（Arithmetical Algebraic Geometry），Purdue University，Harper--Row，New York，1965。
\item[{[2]}] Mumford D.：《病态现象 III》（Pathologies III），Amer. J. Math.，1967，第 94--103 页。Amer. J. Math.，1967，第 94--103 页。
\item[{[3]}] Grothendieck A. 与 Dieudonné J.：EGA，《\Term{algebraic-geometry}{代数几何}基础》（Éléments de Géométrie Algébrique），Pub. Math. Inst. des Hautes Études Sci.，1960。
\item[{[4]}] Grothendieck A.：FGA，《\Term{algebraic-geometry}{代数几何}基础》（Fondements de la géométrie algébrique），摘自 Séminaire Bourbaki，1957--1962。
\item[{[5]}] Grothendieck A. 等：SGA，《\Term{algebraic-geometry}{代数几何}讨论班》（Séminaire de géométrie algébrique），1960--1966，I.H.E.S.。
\item[{[6]}] Jussila O.：X，《\Term{proper}{固有}\Term{algebraic-scheme}{代数概形}上的相\Term{intersection-form}{交形式}体系》（Formalisme des intersections sur les schémas algébriques propres），本讨论班第 X 讲。
\item[{[7]}] Raynaud M.：XII，《\Term{picard-functor}{Picard 函子}的相对\Term{representability}{可表示性}\Term{theorem}{定理}》（Un théorème de représentabilité relative sur le foncteur de Picard），本讨论班第 XII 讲。
\end{enumerate}
